\section{Mô hình hoạch định chuyển động}

Trong phần này, chúng tôi trình bày mô hình toán học của bài toán hoạch định chuyển động cho hệ thống robot tự hành với các ràng buộc vi phân liên tục. Mục tiêu là xây dựng một khung lý thuyết vững chắc để phát triển các thuật toán lập kế hoạch hiệu quả, đồng thời đảm bảo tính khả thi về mặt động lực học và tối ưu hóa hiệu suất vận hành trong môi trường phức tạp chứa vật cản. 

\subsection{Mô hình toán học}

Ta xem xét một quỹ đạo mượt trong $\mathbb{R}^d$. Một quỹ đạo có thể được biểu diễn dưới dạng hàm số liên tục:
\begin{equation}
    \mathbf{x}(t): [0, T] \rightarrow \mathbb{R}^d,
\end{equation}
trong đó $T$ là thời gian hoàn thành quỹ đạo. Quỹ đạo này phải thỏa mãn điều kiện khả vi đến bậc $K$, tức là:
\begin{equation}
    \frac{d^k \mathbf{x}(t)}{dt^k} \text{ tồn tại và liên tục với } k = 1, 2, \ldots, K,
\end{equation}
với đạo hàm bậc 0 là chính hàm số $\mathbf{x}(t)$.

Hàm mục tiêu cần được tối ưu hóa có thể là tổng chi phí năng lượng, thời gian hoàn thành, hoặc một hàm chi phí tổng quát khác. Cụ thể, ta định nghĩa hàm mục tiêu dưới dạng tích phân:
\begin{equation}
    \psi(\mathbf{x}(t), T) := a\Phi(T) + \sum_{m=1}^{M}b_m\int_0^T \left\| \frac{d^m \mathbf{x}(t)}{dt^m} \right\|^2 dt,
\end{equation}
trong đó $\Phi(T)$ là hàm chi phí liên quan đến thời gian, $m$ là bậc đạo hàm đại diện cho đặc tính động lực học cần tối ưu hóa (ví dụ: vận tốc, gia tốc), và $a, b_m$ là các hệ số trọng số điều chỉnh tầm quan trọng tương đối của các thành phần trong hàm mục tiêu.

Tổng quát, bài toán hoạch định chuyển động có thể được mô tả như sau:
\begin{equation}
\begin{aligned}
    & \min_{\mathbf{x}(t), T} & & \psi(\mathbf{x}(t), T) \\
    & \text{subject to} & & \mathbf{x}(0) = \mathbf{x}_{\text{start}}, \quad \mathbf{x}(T) = \mathbf{x}_{\text{goal}}, \\
    & & & \mathbf{x}(t) \in \mathcal{C}_{\text{free}}, \quad \forall t \in [0, T], \\
    & & & \frac{d^k \mathbf{x}(t)}{dt^k} \in \mathcal{D}_k, \quad k = 1, 2, \ldots, K, \quad \forall t \in [0, T], \\
    & & & T_{\min} \leq T \leq T_{\max}.
\end{aligned}
\end{equation}

Trong đó, ràng buộc đầu và cuối đảm bảo quỹ đạo bắt đầu và kết thúc tại các vị trí xác định. Ràng buộc không gian trạng thái yêu cầu quỹ đạo phải nằm trong không gian tự do $\mathcal{C}_{\text{free}}$. Ràng buộc động học giới hạn các đạo hàm bậc $k$ của quỹ đạo trong các tập hợp khả thi $\mathcal{D}_k$, phản ánh các giới hạn vật lý của hệ thống robot. Cuối cùng, ràng buộc về thời gian đảm bảo rằng thời gian hoàn thành quỹ đạo nằm trong khoảng cho phép $[T_{\min}, T_{\max}]$.

Quan sát hàm mục tiêu và ràng buộc, khi tổng thời gian thực hiện quỹ đạo $T$ được coi là một biến quyết định (decision variable) để tối ưu hóa hiệu suất, cấu trúc của hàm mục tiêu tích phân $\int_0^T \| \mathbf{x}^{(m)}(t) \|^2 dt$ trở nên cực kỳ phức tạp. Do giới hạn tích phân thay đổi theo $T$, mối liên hệ giữa các biến trạng thái và thời gian bị ràng buộc một cách phi tuyến tính. 

\subsection{Tham số hóa quỹ đạo dựa trên tọa độ đường đi}

Để giải quyết sự phức tạp của bài toán tối ưu hóa đồng thời hình dạng quỹ đạo và phân bổ thời gian, một phương pháp tiếp cận hiệu quả là thực hiện tham số hóa quỹ đạo thông qua một tọa độ đường đi vô hướng $s \in [0, 1]$ (~\cite{BobrowDubowskyGibson1985TimeOptimal,ShinMcKay1985MinimumTime}). Việc sử dụng biến độc lập $s$ thay cho biến thời gian $t$ cho phép phân rã quỹ đạo thành hai thành phần biệt lập: đặc tính hình học không gian và quy luật tiến triển theo thời gian~\cite{Verscheure2009TimeOptimalPathTracking}. Trong cấu trúc này, quỹ đạo $\mathbf{x}(t)$ được tái cấu trúc thông qua cặp hàm số bao gồm hàm hình dạng $q(s): [0, 1] \rightarrow \mathbb{R}^d$ xác định hình dáng hình học trong không gian cấu hình, và hàm thời gian $h(s): [0, 1] \rightarrow \mathbb{R}$ xác định thời điểm mà robot đi qua vị trí tương ứng với tọa độ $s$. Khi đó, quỹ đạo thực tế được biểu diễn dưới dạng hàm hợp $\mathbf{x}(t) = q(h^{-1}(t))$, trong đó mối liên hệ giữa thời gian và tọa độ đường đi được thiết lập bởi phương trình $t = h(s)$.

Việc tách biệt này mang lại lợi ích đáng kể trong việc thiết kế các quỹ đạo mượt mà bằng cách áp đặt các điều kiện về tính khả vi lên hàm $q(s)$ và $h(s)$ một cách độc lập. Tuy nhiên, phương pháp tham số hóa này làm nảy sinh những thách thức nhất định khi thiết lập các ràng buộc động học cấp cao. Trong khi các hàm chi phí lồi và các ràng buộc về vận tốc vẫn duy trì được tính lồi khi biểu diễn qua $q(s)$ và $h(s)$, thì các ràng buộc liên quan đến đạo hàm bậc hai (như gia tốc) hoặc cao hơn thường mất đi tính lồi vốn có \cite{DeitsTedrake2022}. Điều này gây khó khăn cho các thuật toán tối ưu hóa lồi truyền thống trong việc tìm kiếm nghiệm tối ưu toàn cục cho các bài toán có ràng buộc động lực học phức tạp.

Để khắc phục hạn chế về tính phi lồi của các đạo hàm cấp cao trong khung tham số hóa $q(s)$ và $h(s)$, một chiến lược thực tiễn thường được áp dụng là sử dụng kỹ thuật điều chuẩn (regularization). Thay vì giải quyết trực tiếp các ràng buộc phi lồi khắt khe, hệ thống sẽ bổ sung các thành phần phạt (penalty) tỉ lệ thuận với độ lớn của các đạo hàm bậc cao như $q''(s)$ và $h''(s)$ vào hàm mục tiêu. Đồng thời, hàm thời gian $h(s)$ được ràng buộc chặt chẽ để đảm bảo đạo hàm $h'(s)$ không tiệm cận giá trị không, nhằm duy trì tính đơn điệu tăng của thời gian và đảm bảo tính khả nghịch của phép biến đổi \cite{DeitsTedrake2022}. Cách tiếp cận này cho phép xấp xỉ các ràng buộc kinodynamic trong một không gian tìm kiếm ổn định hơn, giúp bộ lập kế hoạch hội tụ nhanh chóng về các quỹ đạo khả thi về mặt vật lý ngay cả trong các bài toán có số chiều lớn.
